\subsection{Secondary threshold-method sensitivity: ELR/ESR}
The primary $10^{-7}$ constant-current threshold and DIBL results
above are retained unchanged. As an \emph{exploratory, model-dependent
alternative}, we extracted an apparent threshold from the measured
transfer-curve shapes using the linear extrapolation (ELR) and
saturation extrapolation (ESR) conventions.

At $V_{DS}=0.1$~V, a local straight-line fit
$I_D \simeq aV_{GS}+b$ was selected at the maximum
positive slope of the fitted current curve. The conventional
linear-region correction gives
\[
V_{\mathrm{TH,ELR}}=-\frac{b}{a}-\frac{0.1}{2}.
\]
At $V_{DS}=1.2$~V, an independently selected local fit to
$\sqrt{|I_D|}\simeq cV_{GS}+d$ gives the
\emph{saturation-model extrapolate}
\[
V_{\mathrm{TH,ESR}}=-\frac{d}{c},\qquad
\mathrm{DIBL}_{\mathrm{ELR/ESR,app}}=
1000\,\frac{V_{\mathrm{TH,ELR}}-V_{\mathrm{TH,ESR}}}
{1.2-0.1}\quad(\mathrm{mV/V}).
\]
The fits use five neighbouring measured gate-bias points
($\Delta V_{GS}=0.05$~V) and maximize the fitted slope
over all fully contained five-point windows. No curve is
extended or simulated again; a tangent intercept may lie
outside the measured gate interval. We also vary the fitting
window over 3, 5, 7 and 9 points to assess numerical stability.
The raw source is \path{data/part4_idvg_all.csv};
the generated values are in
\path{data/part4_elr_esr_sensitivity.csv}, produced by
\path{tools/gm_dibl_exploratory.py}.

\begin{table}[H]\centering\small
\caption{Alternative, \emph{apparent} thresholds and DIBL for
the four Part 4 samples. These are fit extrapolations under
different low-/high-drain transport approximations, not
direct constant-current threshold crossings.}
\begin{tabular}{@{}rrrrr@{}}\toprule
$N_{\mathrm{CH}}$ (cm$^{-3}$) &
$V_{\mathrm{TH,ELR}}$ (V) &
$V_{\mathrm{TH,ESR}}$ (V) &
Apparent DIBL & 3--9-point range\\
&&& (mV/V) & (mV/V)\\\midrule
$3\times10^{15}$ & 0.1899 & $-0.3574$ & 497.5 & 497.5--497.8\\
$3\times10^{16}$ & 0.2065 & $-0.3296$ & 487.4 & 487.0--487.5\\
$3\times10^{17}$ & 0.2741 & $-0.1822$ & 414.8 & 414.6--417.9\\
$3\times10^{18}$ & 0.6279 & 0.4766 & 137.5 & 136.1--138.9\\
\bottomrule\end{tabular}
\end{table}

\begin{figure}[H]\centering
\begin{tikzpicture}
\begin{semilogxaxis}[
width=.90\linewidth,height=6.3cm,
xlabel={$N_{\mathrm{CH}}$ (cm$^{-3}$)},
ylabel={Extracted drain-bias shift (mV/V)},
xmin=1.5e15,xmax=5e18,ymin=0,ymax=550,
xtick={3e15,3e16,3e17,3e18},
xticklabels={$3\times10^{15}$,$3\times10^{16}$,$3\times10^{17}$,$3\times10^{18}$},
grid=both,
legend style={font=\scriptsize,at={(.03,.04)},anchor=south west}
]
\addplot+[mark=square*,thick]
table[col sep=comma,x=case,y=DIBL_apparent_w5_mV_per_V]
{data/part4_elr_esr_sensitivity.csv};
\addlegendentry{Apparent ELR/ESR shift}
\addplot+[only marks,mark=*,mark size=2.8pt]
coordinates {(3e18,93.57225)};
\addlegendentry{Original $10^{-7}$ criterion}
\end{semilogxaxis}
\end{tikzpicture}
\caption{Doping dependence of the alternative
model-derived gate-voltage shift. The isolated strict-current
DIBL value is plotted for comparison; the two methods
have different definitions and must not be pooled.}
\end{figure}

Both the high-current switching marker and the ELR/ESR
calculation independently suggest that the low-doping
devices are strongly drain-bias sensitive. The ELR/ESR
results are stable to the fitted window length; all
five-point local regressions have $R^2>0.9997$.
However, high local linearity is \emph{not} proof that
a short-channel, high-leakage structure obeys the
ideal linear-region or square-law saturation expressions.
Indeed the $3\times10^{18}$~cm$^{-3}$ device gives
137.5~mV/V by ELR/ESR but 93.57~mV/V by the
original constant-current criterion. Consequently
these numbers are reported only as
\emph{apparent, extraction-method-dependent shifts};
they do not retroactively establish the missing strict-current
thresholds and cannot by themselves prove quantitative
drain-induced source-barrier lowering.
